\subsubsection{唯一性}\label{subsec:conservation-uniqueness}

现已讨论了几种选择准则, 并指出它们之间的一些关系. 原先的目标是得到某种唯一性结果. 做过这些工作之后, 是否已能做到这一点? 一般而言, 答案是"不能". 虽然所提出的准则排除了最明显的"物理上不合理"的弱解, 但一般的存在性与唯一性问题仍是未解决的. 在撰写本书时, 这是一个十分活跃的研究领域. 下文概述一些特殊情形的结果.

对具有强凸 $f$ 的标量守恒律, 存在性与唯一性问题已基本解决. 对真正非线性的方程组, 已知对全变差较小的初始数据存在解 (但尚未知唯一性). 对 $p$ 方程组, 假设强凸性时, 已知结果多得多. 任意初始数据都存在解, 且已在分片光滑解类中证明了唯一性. 相关论述参见~\cite{Sm}.

对其他方程组, 也有许多专门的结果, 例如以特定方式违背真正非线性的方程组, 以及气体动力学方程组. 存在性结果通常基于对近似解建立估计, 再提取收敛子列. 这类近似解通常来自有限差分格式, 或者来自向方程中添加"黏性"项. 这一领域的一些主要贡献者是 Lax, Glimm, DiPerna, Tartar, Godunov, Liu, Smoller 和 Oleinik. 尽管作出了所有这些努力, 这一领域的一般答案仍难以获得. 事实上, 近来已有反例表明, 通常的容许性条件不能保证唯一性~\cite{Se}. 当然, 现实问题通常具有不止一个空间维度. 在这种情形下, 几乎所有问题都尚未解决.

\paragraph*{Problems}

\begin{enumerate}[label=\thesection.\arabic*.,ref=\thesection.\arabic*,start=1,series=book-problems]

\item\label{ex:p-system-hyperbolicity} 证明, 若 $p'<0$, 则 $p$ 方程组是双曲型的.

\item\label{ex:gas-dynamics-hyperbolicity} 给出本构函数应满足的条件, 以保证两个气体动力学方程组是双曲型的.

\item\label{ex:conservation-characteristics-proof} 证明 Lemma~\ref{lem:conservation-characteristics}.

\item\label{ex:riccati-blowup-proof} 证明 Lemma~\ref{lem:riccati-finite-time-blowup}.

\item\label{ex:p-system-blowup-proof} 证明 Theorem~\ref{thm:p-system-finite-time-blowup}.

\item\label{ex:sketch-mixed-rarefaction-shock} 对 \eqref{eq:burgers-continuum-solution-data} 给出的解集, 画出特征线图与波前.

\item\label{ex:convexity-lax-scalar-shock} 设 $f$ 凸, $u$ 是 \eqref{eq:scalar-conservation-law-quasilinear} 的分片光滑弱解, 具有有限个跳跃. 证明, 若 $u$ 作为 $x$ 的函数单调递减, 则它在每个间断处都满足 Lax 激波条件. 举例说明, 若 $f$ 非凸, 则这一结论不成立.

\item\label{ex:entropy-jump-condition-proof} 证明 Lemma~\ref{lem:entropy-jump-condition}.

\item\label{ex:riemann-invariant-gradient-basis} 证明, \eqref{eq:rarefaction-invariant-gradient-basis} 中定义的集合 $V$ 是 $\mathbb{R}^n$ 的一组基. Hint: 当 $i\ne j$ 时, $\bm{r}_i$ 与 $\bm{l}_j$ 有什么关系?

\item\label{ex:p-system-shock-reparametrization} 证明, \eqref{eq:p-system-admissible-first-shock-states} 和 \eqref{eq:p-system-admissible-second-shock-states} 中定义的曲线所给出的状态, 与对应的激波速度一起, 分别满足 \eqref{eq:p-system-back-shock-speeds} 和 \eqref{eq:p-system-front-shock-speeds}. Hint: 在开平方之前利用 $p$ 的凸性.

\item\label{ex:p-system-fast-family-covering} 证明, 快族单值覆盖 $\bar{\bm{u}}$ 的一个邻域.

\item\label{ex:eulerian-lagrangian-gas-equivalence} 证明, 对光滑解, Eulerian 与 Lagrangian 气体动力学等价. 对弱解会出现什么困难?

\end{enumerate}
